\section{Predication as the seed out of potential}

\begin{proposition}[Entity and predicate]\label{thm:entity-predicate}\lean{Obsmetrologia.Entity}\leanok\uses{def:token, def:claim}
An \emph{entity} is a token: a mark in identity-carrying role, the
subject/object position of a predication. An \emph{entity} is
$\texttt{Entity} := \texttt{Token}$ and a \emph{predicate} is
$\texttt{Predicate} := \texttt{Claim}$ --- abbreviations of the two typed
hands of \cref{def:token,def:claim}, not of the bare mark underneath them.
Typing proper enters further downstream, at positing and valuation.
\end{proposition}

\begin{proposition}[The predication cell]\label{thm:predication}\lean{Obsmetrologia.PredicationCell}\leanok\uses{thm:entity-predicate, thm:articulation-cell-corner}
The predication cell is the oriented triple $(s,p,o) : \{\text{subject} :
\texttt{Token},\ \text{predicate} : \texttt{Claim},\ \text{object} :
\texttt{Token}\}$: subject entity, predicate, and object entity. Orientation
belongs to the datum, so the same predicate read from opposite endpoints is a
different cell off the diagonal.
\end{proposition}

\begin{proposition}[The predication cell as the point value is carried at]\label{thm:claim-point}\lean{Obsmetrologia.PredicationCell}\leanok\uses{thm:predication}
A claim is the utterance's right-hand carrier, not the cell. The predication
cell is the generated point at which that claim's later value, settlement, and
descent can be carried. It is not yet a site object, a cover, or a glued
documentation carrier; those structures are generated later.
\end{proposition}

\begin{proposition}[The reflexive seed]\label{thm:predication-seed}\lean{Obsmetrologia.PredicationCell.seed}\leanok\uses{thm:predication, def:utterance}
The seed of an utterance with token $t$ and claim $c$ is the reflexive
predication cell $(t,c,t)$. Predicating first creates the diagonal assertion
position before later constructions move off the diagonal.
\end{proposition}

\begin{theorem}[The claim seed records the claim]\label{thm:claim-seed-predicate}\lean{Obsmetrologia.PredicationCell.seed_predicate}\leanok\uses{thm:claim-point, thm:predication-seed}
The predication cell generated by an utterance's seed keeps the utterance's
claim as its predicate coordinate. The claim itself is that right-hand carrier
(\texttt{Utterance.claim}); the cell only carries it, which is why the cell is
named \emph{PredicationCell} rather than \emph{Claim}.
\end{theorem}

\begin{proposition}[Predicating proper]\label{thm:predicating-proper}\lean{Obsmetrologia.PredicationCell.predicating}\leanok\uses{thm:predication-seed, thm:articulation-cell-initial}
Predicating is defined on an articulation cell only with a proof that the
cell is in the potential corner. Under that incidence condition it returns
the seed of the cell's utterance.
\end{proposition}
